\section{Introduction}

Many scientific systems are governed by partial differential equations whose large-scale structure is known while crucial constitutive closures remain unavailable. This setting is common in transport, pattern formation, and continuum modeling: the conservation law or balance law is understood, but the diffusion, mobility, reaction, or source terms must still be inferred from data. In such problems, the relevant scientific object is not only a predictive simulator but the closure law itself.

This paper studies closure discovery in nonlinear reaction-diffusion systems of the form
\begin{equation}
u_t = \partial_x \left( D(u) u_x \right) + R(u),
\label{eq:rd_main}
\end{equation}
where the PDE structure is known but the constitutive relations $D(u)$ and $R(u)$ are unknown. Given spatiotemporal observations of $u(x,t)$, we seek to recover these closures in a form that is numerically valid, scientifically interpretable, and verifiable by forward simulation.

The main challenge is that closure discovery is an inverse problem rather than a pure prediction problem. A small residual on observed trajectories, or even a good short-horizon rollout, does not by itself imply correct constitutive recovery. Under limited excitation, different closure pairs may compensate each other on the observed support and produce similar trajectory-level behavior. This creates a gap between fitting the observed dynamics and identifying the true law. Closing that gap is the main methodological issue addressed in this work.

This perspective immediately changes how methods should be evaluated. In matched-library settings, where the true closure lies in a low-order polynomial dictionary, weak polynomial baselines can behave like correctly specified reference estimators and should be treated as strong comparators rather than weak strawmen. In mismatch settings, where the true closure lies outside that restricted library, flexible neural surrogates become relevant as numerical constitutive approximators. However, even then, a neural surrogate is not yet a scientific law. It still must be compressed into a symbolic form and then reinserted into the PDE to check whether the compressed law preserves the learned dynamics.

Our paper therefore adopts a three-stage neural-symbolic view of closure discovery. Stage~1 learns numerical constitutive surrogates in weak form. Stage~2 compresses those surrogates into restricted symbolic families. Stage~3 validates the symbolic laws by forward re-simulation on unseen initial conditions. This decomposition lets us ask a sharper question than ``does the method get low loss?'': what part of the pipeline is actually responsible for constitutive error?

The answer supported by our theory and experiments is clear. Symbolic compression can preserve learned constitutive behavior with very small additional approximation error and almost no rollout degradation, but it does not automatically repair constitutive bias inherited from the neural surrogate. Across clean, noisy, and sparse observation settings, the symbolic true closure error closely tracks the neural true closure error, and the bias inheritance ratio remains near one. The main bottleneck in neural-symbolic closure discovery is therefore Stage~1 numerical constitutive recovery rather than Stage~2 symbolic compression.

The contributions of this paper are fourfold.
\begin{enumerate}[leftmargin=1.5em]
\item We formulate closure discovery under known PDE structure as a weak-form inverse problem with explicit attention to identifiability, excitation coverage, and function-class mismatch.
\item We propose a neural-symbolic pipeline that separates numerical constitutive recovery, restricted symbolic compression, and forward validation by PDE re-simulation.
\item We provide a theory spine explaining why low residual does not guarantee true closure recovery, why matched-library baselines can be unusually strong in the correctly specified regime, and why symbolic compression preserves rather than repairs constitutive bias.
\item We benchmark the method against strong and weak polynomial baselines across clean, noisy, and sparse regimes, showing that the central mechanism is bias preservation through compression rather than symbolic correction.
\end{enumerate}
